\documentclass[11pt]{article}
\usepackage{fontspec}
\setmainfont{TeX Gyre Pagella}
\setsansfont{TeX Gyre Heros}[Scale=MatchLowercase]
\setmonofont{Latin Modern Mono}
\newfontfamily\termes{TeX Gyre Termes}
\usepackage{xcolor}
\usepackage[colorlinks]{hyperref}
\usepackage{booktabs}
\title{XeTeX in partex}\author{Somebody}\date{1 January 2026}
\begin{document}
\maketitle
\tableofcontents
\section{Fonts}\label{sec:a}
Pagella text with ligatures fi fl and figures 1234567890. \textsf{Heros sans}
and \texttt{mono text}. \textit{Italic} \textbf{bold} \textsc{Small Caps}.
{\addfontfeatures{Numbers=OldStyle}0123456789} --- en–dash, a footnote\footnote{The
note — with “quotes”.} and a reference to Section~\ref{sec:b} on page~\pageref{sec:b}.
\begin{itemize}
\item First item, \emph{emphasised} and \textbf{bold}.
\item Second \textcolor{red}{red} item.
\end{itemize}
\section{Hyphenation}\label{sec:b}
\parbox{4cm}{\termes Characteristically internationalization encyclopedia
``hyphenation'' (parenthetical) incomprehensibilities, uncharacteristically
representational extraordinarily; co-operation—dashes—everywhere.}
\showhyphens{internationalization incomprehensibilities}
\begin{table}[h]\centering
\begin{tabular}{lrc}\toprule Name & Value & Unit\\\midrule Alpha & 1.5 & m\\ Beta & 22 & s\\\bottomrule\end{tabular}
\caption{A table}\label{tab:x}
\end{table}
See Table~\ref{tab:x}. Ünïcödé: ÀÉÎÕÜ àéîõü ß Œœ Ææ ł.
\end{document}
